\documentclass[19pt,pdf,hyperref={unicode}]{beamer}
\beamertemplatenavigationsymbolsempty
\setbeamertemplate{footline}[page number]

\usepackage[T2A]{fontenc}
\usepackage[utf8]{inputenc}
\usepackage[russian]{babel}
\usepackage{bm}
\usepackage{multirow}
\usepackage{ragged2e}
\usepackage{indentfirst}
\usepackage{multicol}
\usepackage{subfig}
\usepackage{amsmath,amssymb}
\usepackage{enumerate}
\usepackage{mathtools}
\usepackage{comment}
\usepackage[all]{xy}
\usepackage{tikz}
\usetikzlibrary{positioning,arrows}
\usepackage{graphicx}
\usepackage{tabularx}
\usepackage{booktabs}

\definecolor{name}{rgb}{0.5,0.5,0.5}

\newcommand{\R}{\mathbb{R}}
\newcommand{\Normal}{\mathcal{N}}
\newcommand{\Cat}{\mathrm{Cat}}
\newcommand{\Loss}{\mathcal{L}}
\newcommand{\doop}{\operatorname{do}}
\newcommand{\Cov}{\operatorname{Cov}}
\newcommand{\diag}{\operatorname{diag}}
\newcommand{\Xobs}{X^{\mathrm{obs}}}
\newcommand{\Yobs}{Y^{\mathrm{obs}}}

\title[Probabilistic Causal CCM]{Вероятностные модели причинно-следственных связей\\в сходящихся отображениях}
\author{Иконников Марк}
\date{май 2026}

\begin{document}

% Slide 1
\begin{frame}
    \titlepage
\end{frame}

% Slide 2
\begin{frame}{Проблема}
Даны два высокоразмерных временных ряда
\[
\mathcal{D} = \{(\Xobs(t), \Yobs(t))\}_{t=1}^{T}, \qquad
\Xobs(t),\Yobs(t) \in \R^d, \quad d \gg 1.
\]
Требуется восстановить интерпретируемую причинную структуру между сигналами, имея только наблюдательные данные.

\begin{center}
\includegraphics[width=0.5\textwidth]{Фазовые траектории_v2.png}
\end{center}
\end{frame}

% Slide 3
\begin{frame}{Цель работы}
\begin{block}{Основная цель}
Построить структурную причинную модель, которая по наблюдательным данным выделяет скрытые режимы и поддерживает интервенционные запросы.
\end{block}

\begin{block}{Идея}
Мы не восстанавливаем граф напрямую в исходном пространстве $\R^d$. Сначала переходим в CCA-латент, строим там SCM, а затем переносим найденные связи обратно.
\end{block}

\begin{figure}
    \centering
    \includegraphics[width=0.7\textwidth]{Снимок экрана 2026-05-23 142331.png}
    \caption{Classic CCA}
    \label{fig:cca-classic}
\end{figure}
\end{frame}

% Slide 4
\begin{frame}{Ключевая идея подхода}
\begin{columns}
\begin{column}{0.55\textwidth}
\begin{block}{Latent SCM}
Работаем с минимальным графом
\[
    Z \to X, \qquad Z \to Y, \qquad X \to Y.
\]
Скрытый $Z$ отвечает за общий фактор, а канал $X \to Y$ --- за прямое влияние между сигналами.
\end{block}

\begin{block}{Обучаемая интервенция}
Оператор $\doop(Z=z_0)$ заменяется обучаемой стохастической политикой-смесью. Компоненты смеси соответствуют ненаблюдаемым режимам данных.
\end{block}
\end{column}

\begin{column}{0.45\textwidth}
\centering
\begin{tikzpicture}[node distance=1.8cm, thick, >=stealth]
    \node[circle, draw, fill=blue!10, minimum size=0.9cm] (Z) {$Z$};
    \node[circle, draw, fill=blue!10, below left of=Z, minimum size=0.9cm] (X) {$X$};
    \node[circle, draw, fill=blue!10, below right of=Z, minimum size=0.9cm] (Y) {$Y$};
    \draw[->] (Z) -- (X);
    \draw[->] (Z) -- (Y);
    \draw[->] (X) -- (Y);
\end{tikzpicture}

\vspace{0.5cm}
\fbox{\parbox{0.9\linewidth}{\centering Observed $X,Y$ $\to$ CCA latent $\to$ SCM $\to$ graph}}
\end{column}
\end{columns}
\end{frame}

% Slide 5
\begin{frame}{Структурная модель в латенте}
Фиксируем линейно-гауссову SCM:
\[
\begin{aligned}
    Z &\sim \pi_\xi(Z), \\
    X &= A Z + \varepsilon_x, \qquad \varepsilon_x \sim \Normal(0, \Sigma_x), \\
    Y &= B Z + C X + \varepsilon_y, \qquad \varepsilon_y \sim \Normal(0, \Sigma_y).
\end{aligned}
\]

\begin{block}{Смысл параметров}
$A$ описывает канал $Z \to X$, $B$ --- канал $Z \to Y$, а $C$ --- прямой канал $X \to Y$. Все три матрицы являются общими для всех режимов.
\end{block}

\begin{block}{Каузальная интерпретация}
Модель явно разделяет общий источник зависимости и прямое влияние, поэтому интервенция на $Z$ или $X$ имеет смысл внутри заданной SCM.
\end{block}
\end{frame}

% Slide 6
\begin{frame}{Обучаемый do-оператор}
Задаем распределение интервенций как смесь гауссиан:
\[
    \pi_\xi(Z) = \sum_{e=1}^{E_{\max}} w_e \Normal(Z;\mu_e,\Sigma_e),
    \qquad \sum_e w_e = 1, \quad w_e \ge 0.
\]
Каждая компонента задает soft-интервенцию
\[
    \doop_e(Z) \equiv \doop\big(Z \sim \Normal(\mu_e,\Sigma_e)\big).
\]

\begin{block}{Предельный случай}
При $\Sigma_e \to 0$ компонента смеси вырождается в атомарную интервенцию $\doop(Z=\mu_e)$.
\end{block}
\end{frame}

% Slide 7
\begin{frame}{3CA Algorithm}
\begin{enumerate}
    \item \textbf{Выбор лага:} ищем $\tau^*$ через максимум старшего CCA-коэффициента.
    \item \textbf{CCA-сжатие:} переходим от $\R^d$ к $\R^k$.
    \item \textbf{Causal Mixture Model:} обучаем SCM с общими параметрами $\theta=(A,B,C)$ и смесью режимов.
    \item \textbf{EM-обучение:} выводим responsibilities $r_{t,e}$ и параметры режимов.
    \item \textbf{Извлечение графа:} получаем per-regime оператор $F'_e$ и переносим его в исходное пространство.
\end{enumerate}
\end{frame}

% Slide 8
\begin{frame}{Генеративная модель, loss и EM}
\small
\begin{block}{Causal mixture}
\[
 p(X,Y,Z,e\mid\theta,\xi)=p(e\mid\xi)p(Z\mid e;\xi)p(X\mid Z;\theta)p(Y\mid X,Z;\theta).
\]
\[
 p(X_t,Y_t\mid\theta,\xi)=\sum_{e=1}^{E_{\max}} w_e
 \int p(X_t\mid Z;\theta)p(Y_t\mid X_t,Z;\theta)\Normal(Z;\mu_e,\Sigma_e)dZ.
\]
\end{block}

\begin{block}{Полученный $\mathcal{L}$}
\[
\Loss(\theta,\xi)=
-\sum_{t=1}^{T}\log\sum_{e=1}^{E_{\max}}w_e p(X_t,Y_t\mid\doop_e;\theta)
-\log p(\theta)-\log p(w).
\]
\end{block}

\begin{block}{EM-обучение}
\[
 r_{t,e}=\frac{w_e p(X_t,Y_t\mid e;\theta,\xi)}
 {\sum_{e'}w_{e'}p(X_t,Y_t\mid e';\theta,\xi)}.
\]
На M-шаге $\theta=(A,B,C)$ обновляются общими weighted normal equations, а $w_e,\mu_e,\Sigma_e$ --- стандартными формулами GMM.
\end{block}
\end{frame}

% Slide 9
\begin{frame}{Извлечение каузального графа}
После обучения для каждого режима строим оператор связи $X \to Y$ в латенте:
\[
    F'_e = \hat C + \hat B\,\Cov_e(Z,X)\,\Cov_e(X,X)^{-1}.
\]
Далее используем SVD-разложение
\[
    F'_e = U'_e \Lambda'_e {V'_e}^{\top},
    \qquad
    \Lambda'_e = \diag(\lambda'_{e,1},\dots,\lambda'_{e,k}).
\]

\begin{block}{Восстановление  графа}
CCA дает линейные декодеры $D_X,D_Y$. Через них переносим найденные структурные параметры:
\[
    A^{\mathrm{obs}} = D_X\hat A, \qquad
    B^{\mathrm{obs}} = D_Y\hat B.
\]
Оператор связи в исходных координатах:
\[
    F^{\mathrm{obs}}_e = D_Y F'_e D_X^+.
\]

\end{block}
\end{frame}

% Slide 11
\begin{frame}{Теоретические утверждения}
\small
\begin{block}{Утверждение 1: идентифицируемость SCM в латенте}
Пусть SCM линейно-гауссова, истинное число режимов ограничено, а моменты компонент смеси различимы.

Тогда глобальный максимум $\Loss$ восстанавливает структурные параметры $(A,B,C)$ с точностью до перестановки и смены знака латентных компонент.
\end{block}

\begin{block}{Теорема 2: корректность переноса SCM в исходное пространство}
Пусть $\mathrm{SCM}_{\mathrm{latent}}$ --- обученная SCM с параметрами $(\hat A,\hat B,\hat C)$ в латенте, а $D_X,D_Y$ --- CCA-декодеры. Тогда индуцированная SCM в $\R^d$ с
\[
A^{\mathrm{obs}}=D_X\hat A,\quad B^{\mathrm{obs}}=D_Y\hat B,\quad F_e^{\mathrm{obs}}=D_YF'_eD_X^+
\]
сохраняет интервенционную семантику:
\[
 p^{\mathrm{obs}}(\Yobs\mid\doop(\Xobs))=
 p^{\mathrm{latent}}(Y\mid\doop(X))\Big|_{X=D_X^+\Xobs,\;Y=D_Y^+\Yobs}.
\]
\end{block}
\end{frame}

% Slide 12
\begin{frame}{Связь с CCM}
\begin{block}{Классический CCM}
CCM проверяет, можно ли по shadow manifold одного ряда восстановить состояние другого. Рост качества cross-mapping при увеличении длины библиотеки интерпретируется как признак причинной связи.
\end{block}

\begin{block}{Наша модификация}
Мы добавляем к CCM латентную интервенционную модель. CCA выделяет совместный низкоразмерный сигнал, mixture-модель разделяет режимы, а per-regime операторы $F'_e$ дают слой, внутри которого можно применять CCM или сравнивать causal strength.
\end{block}
\end{frame}

% Slide 13
\begin{frame}{Метрики качества}
\small
\begin{tabularx}{\textwidth}{lX}
\toprule
\textbf{Метрика} & \textbf{Что измеряет} \\
\midrule
Predictive error & Насколько хорошо модель предсказывает $Y_t$ по $X_{t-\tau}$ \\
Cross-map skill & Растет ли качество CCM при увеличении library length \\
Regime stability & Устойчивы ли найденные режимы при разных инициализациях EM \\
Graph sparsity & Получается ли разреженный и интерпретируемый оператор $F'_e$ \\
Intervention consistency & Согласованы ли ответы модели на $\doop_e(Z)$ между режимами \\
\bottomrule
\end{tabularx}

\vspace{0.3cm}
\begin{block}{Базовые сравнения}
Сравниваем с CCA без каузального слоя, с единственным режимом $E=1$, и с обычным CCM без латентного разделения режимов.
\end{block}
\end{frame}

% Slide 14
\begin{frame}{Экспериментальный сетап}
\begin{block}{Данные}
Парные многомерные временные ряды с датчиков движения: акселерометр, гироскоп и ориентация.
\end{block}

\begin{figure}
    \centering
    \includegraphics[width=1.0\textwidth]{all_series.png}
    \caption{Временные ряды}
    \label{fig:all-series}
\end{figure}
\end{frame}

% Slide 15
\begin{frame}{Результаты}
\begin{figure}
    \centering
    \includegraphics[width=0.98\textwidth]{acc_vs_dim_v2.png}
    \label{fig:acc-vs-dim}
\end{figure}

\vspace{-0.2cm}
\begin{figure}
    \centering
    \includegraphics[width=0.78\textwidth]{regimes.png}
    \label{fig:regimes}
\end{figure}
\end{frame}

% Slide 16
\begin{frame}{Результаты: модель и предсказание}
\begin{columns}[T]
\begin{column}{0.48\textwidth}
    \centering
    \includegraphics[width=0.95\textwidth]{3ca.png}
\end{column}
\begin{column}{0.48\textwidth}
    \centering
    \includegraphics[width=0.95\textwidth]{causal_graph.png}
\end{column}
\end{columns}

\vspace{0.35cm}
\begin{center}
    \includegraphics[width=0.82\textwidth]{ci_cca_pred.png}
\end{center}
\end{frame}

% Slide 18
\begin{frame}{Тезисы, выносимые на защиту}
\begin{enumerate}
    \item \textbf{Вид $\doop$-интервенции.} Оператор $\doop(Z=z_0)$ формализуется как обучаемая стохастическая смесь
    \[
        \pi_\xi(Z)=\sum_e w_e\Normal(Z;\mu_e,\Sigma_e),
    \]
    где компоненты соответствуют скрытым режимам данных.
    \item \textbf{Алгоритм.} Предложен 3CA-пайплайн: выбор лага, CCA-сжатие, causal mixture SCM, EM-обучение, извлечение и перенос графа.
    \item \textbf{Теорема об идентифицируемости.} При различимых режимах и инвариантных $(A,B,C)$ структурные параметры восстанавливаются в латенте с точностью до стандартных неоднозначностей.
    \item \textbf{Теорема о переносе.} CCA-декодеры индуцируют SCM в исходном пространстве и сохраняют интервенционную семантику.
\end{enumerate}
\end{frame}

% Slide 19
\begin{frame}{Литература}
\tiny
\begin{tabularx}{\textwidth}{l c X}
\toprule
\textbf{Авторы} & \textbf{Год} & \textbf{Публикация} \\
\midrule
Klami, Virtanen, Kaski & 2013 & Journal of Machine Learning Research, 14:965--1003 \\
McLachlan, Peel & 2000 & Wiley \\
Dempster, Laird, Rubin & 1977 & Journal of the Royal Statistical Society B, 39(1):1--38 \\
Correa, Bareinboim & 2020 & AAAI \\
Khemakhem, Kingma, Monti, Hyv\"arinen & 2020 & AISTATS \\
Sch\"olkopf et al. & 2021 & Proceedings of the IEEE, 109(5):612--634 \\
Beal & 2003 & PhD thesis, University College London \\
Teicher & 1963 & Annals of Mathematical Statistics, 34:1265--1269 \\
Rubenstein et al. & 2017 & UAI \\
Pearl & 2009 & Cambridge University Press, 2nd edition \\
Sugihara et al. & 2012 & Science, 338(6106):496--500 \\
\bottomrule
\end{tabularx}
\end{frame}

\end{document}
